\documentclass[amssymb,pra,12pt,aps,notitlepage]{revtex4-1}
\usepackage{mathptmx,amsmath, multirow, graphicx}
\usepackage{parskip}

\begin{document}

\title{Putnam POTD November 24 2025}
\maketitle

\section{Problem}
Let $D_n$ denote the value of the $(n-1) \times (n-1)$ determinant
\[
\begin{vmatrix}
3 & 1 & 1 & 1 & \cdots & 1 \\
1 & 4 & 1 & 1 & \cdots & 1 \\
1 & 1 & 5 & 1 & \cdots & 1 \\
1 & 1 & 1 & 6 & \cdots & 1 \\
\vdots & \vdots & \vdots & \vdots & \ddots & \vdots \\
1 & 1 & 1 & 1 & \cdots & n+1
\end{vmatrix}.
\]
Is the set $\{D_n/n!\}_{n \ge 2}$ bounded?

\section{Solution}
First, subtract the \textbf{first row from the last row} to get:
\[
D_n = \begin{vmatrix}
3 & 1 & 1 & 1 & \cdots & 1 \\
1 & 4 & 1 & 1 & \cdots & 1 \\
1 & 1 & 5 & 1 & \cdots & 1 \\
\vdots & \vdots & \vdots & \vdots & \ddots & \vdots \\
1 & 1 & 1 & 1 & \cdots & n \\
-2 & 0 & 0 & 0 & \cdots & n
\end{vmatrix}.
\]
Expanding along the last row gives:
\[
D_n = nD_{n-1} + (-2)(-1)^{(n-1)+1} A_{n-1} = nD_{n-1} + 2(-1)^{n} A_{n-1},
\]
where $A_{n-1}$ is the determinant of the $(n-2) \times (n-2)$ minor obtained by removing the first column and the last row. Explicitly:
\[
A_{n-1} = \begin{vmatrix}
1 & 1 & 1 & \cdots & 1 \\
4 & 1 & 1 & \cdots & 1 \\
1 & 5 & 1 & \cdots & 1 \\
\vdots & \vdots & \ddots & \vdots \\
1 & 1 & \cdots & n & 1
\end{vmatrix}.
\]
Subtract the first row from every other row. The resulting matrix has non-zero entries on the diagonal shifted by one column (e.g., $3, 4, \dots, n-1$) and zeros elsewhere below the first row, except for the last row which becomes $[0, \dots, n-1, 0]$:
\[
A_{n-1} = \begin{vmatrix}
1 & 1 & 1 & \cdots & 1 \\
3 & 0 & 0 & \cdots & 0 \\
0 & 4 & 0 & \cdots & 0 \\
\vdots & \vdots & \ddots & \vdots \\
0 & 0 & \cdots & n-1 & 0
\end{vmatrix}.
\]
Expanding along the last column, the only non-zero contribution comes from the top-right $1$:
\[
A_{n-1} = (-1)^{1+(n-2)} \cdot 1 \cdot \left( \prod_{k=3}^{n} (k-1) \right) = (-1)^{n-1} \frac{(n-1)!}{2}.
\]
Substituting back into the expression for $D_n$:
\[
D_n = nD_{n-1} + 2(-1)^{n} \left[ (-1)^{n-1} \frac{(n-1)!}{2} \right].
\]
Since $(-1)^n (-1)^{n-1} = -1$, this simplifies to:
\[
D_n = nD_{n-1} + (n-1)!.
\]
Let $B_n = D_n/n!$. Dividing the recurrence by $n!$:
\[
\frac{D_n}{n!} = \frac{nD_{n-1}}{n!} + \frac{(n-1)!}{n!} \implies B_n = B_{n-1} + \frac{1}{n}.
\]
For $n=2$, $D_2 = |3| = 3$, so $B_2 = 3/2$. Note that $B_n$ satisfies the recurrence of the Harmonic series $H_n$. Specifically, $B_n = H_n$. Since $\lim_{n \to \infty} H_n = \infty$, the set is unbounded.

\end{document}